\section{System Architecture}\label{sec:system-architecture}
This section describes the challenges encountered during implementation of proposed design and techniques used to solve them.
The design expects full independence between system's components, in practice, some coupling emerges due to accessing shared
resources, execution order, communication, and code reuse.

The architecture follows a hybrid approach combining:
\begin{itemize}
    \item Centralized data access through shared database connections.
    \item Globally accessible Request and Response objects.
    \item Interface-driven communication between components.
    \item Developer constants are defined in respective classes serving static key-value look-up tables.
\end{itemize}

This approach provides easily accessible datasets, that are needed in all parts of the system.

\subsection{Component-Level Architecture}\label{subsec:component-architecture}
This section addresses challenges faced during implementation of each component.



\subsubsection*{Routing}
The routing mechanism is responsible for mapping request paths to actions.
During implementation, several challenges arise:

\begin{itemize}
    \item Internal structure - The initial solution used naive tree implementation, which lacked proper edge definitions.
        This problem manifested through convoluted rules during traversal, that made adding more features increasingly difficult.
        The structural problem was solved by proper implementation of general graph collection and using modified graph
        traversal algorithms.
        Implementation of proper graph collection enabled adding data to connections for reasons beyond routing.
        Simplest example of this feature are menu routes that contain route definition, labels for each segment, and optional icon.
    \item Ambiguity of routes - Routes are dynamic templates for request paths including static and variable parts of a path's segment.
        Segments are marked by their type, which holds precedence rules.
        Search results are aggregated and sorted by the precedence rules to ensuring deterministic order of search results.
    \item Route template evaluation - The system may include hundreds of routes and each of them may be evaluated during search.
        Each route is parsed into intermediate object representation by custom-made route compiler using the compiler architecture~\parencite{appel2000_compiler_design}.
\end{itemize}

Examples of routes:
\begin{verbatim}
    / - Get current node.
    /** - Match all segments after **. This definition is
        least descriptive thus it has the lowest prescedence.
    /foo - Match static "foo"
    /[foo] - Save first segment as "foo" variable
    /foo/[bar] - Match first segment as static "foo"
        and save the value of second segment to the "bar" variable
    /foo/bar-[id] - Match first segment as static "foo"
        and match the second segment as "bar-"
        and save the remainder to the "id" variable
\end{verbatim}



\subsubsection*{Database Layer}
Tha main challenge when developing the database abstraction, was providing clear and expressive configuration.
The implementation of the configuration had to be performant enough to be accessed multiple times for each outgoing query.
Working with JDBC (Java Database Connection) library inspired the code with metadata approach using class attributes.
The database layer abstracts tables through PHP classes.
Each model declares its database connection, associated table, and table columns through PHP attributes.

Example of mapping users table to the User class:
\begin{verbatim}
    #[Database(App::DATABASE)]
    #[Table('users')]
    class User extends Model {
        #[Column('id_user', Column::TYPE_INTEGER, primaryKey: true)]
        protected int $id;

        #[Column('name', Column::TYPE_STRING)]
        protected string $name;
    }
\end{verbatim}

From this definition, the system is capable of:

\begin{itemize}
    \item Executing standard fetch queries (first, all, fromId, and count number of records for given condition).
    \item Creating new User records in database.
    \item Updating a User record.
    \item Deleting a User record.
\end{itemize}

All Query function uses parametrized query if needed.
For example to remove record from users table with id 4.

\begin{verbatim}
    // Creates DeleteQuery object
    $query = Sql::delete('users');

    // DELETE query is requred to have where clause
    // Query::infer automatically infers the parameter type (TYPE_INTEGER)
    $query->where(Query::infer('id_user = ?', 4));
\end{verbatim}

If the developers use provided Query classes correctly there shall never be SQL injection.
The Database Layer provides default fetching functions first and all which uses the same Query object.
Default implementation of first and all functions return one or many instances of model objects.
Fetch function uses the projection parameter to add developer defined SQL query projection optimization.
Implementation of Database Layer does not use the \("\)SELECT *\("\) projection variant although SelectQuery objects
uses it if projection is not set.
Not defining projection for SelectQuery is not encouraged.

\begin{verbatim}
    // Returns record with the id 4
    // populating all column properties on the User object
    // SELECT id_user, name FROM ...
    User::first(where: Query::infer('id_user = ?', 4));

    // Returns record with the id 4
    // populating only columns defined in projection array
    // SELECT name FROM ...
    User::first(
        projection: ['name'],
        where: Query::infer('id_user = ?', 4)
    );

    // When using select query directly projection
    // may not be set and thus the * is used.
    // Uses SELECT * FROM ...
    $query = Sql::select('users')
        ->where(Query::infer('id_user = ?', 4));
\end{verbatim}

Defining custom ORM database interface provides control how the features are implemented and integrated into the system.



\subsubsection*{Asset Loading (SideLoader)}
The SideLoader component introduces asynchronous asset handling, which complicates an otherwise synchronous server-side
rendering pipeline.

Observed issues:
\begin{itemize}
    \item Context independent importing - The core problem that SideLoader addresses is importing assets for code that
        is injected into the HTML page during rendering or loaded asynchronously when HTML components are fetched through AJAX\@.
        The solution uses unique identifier for each asset file and endpoint that accepts list of identifiers and returns asset blob.
        During rendering all assets are caught by the corresponding import functions and list of identifiers is generated.
        If response includes HTML head element, the import call to the endpoint is injected into the head otherwise
        the list is passed through custom HTTP header in the response and JavaScript injects import element into the head.
    \item Dependency resolution - For the sake of simplicity, the dependencies are not resolved and all assets need to imported
        explicitly.
    \item Debugging - JavaScript files requires certain order of importing the scripts.
        Static loading for all pages of commonly used libraries solves most of the issues with incorrect import order.
        For the remaining cases, custom events are dispatched when code is imported.
\end{itemize}



\subsubsection*{Translation System}
The translation system uses static text as keys, which simplifies integration but introduces constraints:

\begin{itemize}
    \item Coupling - If default language is changed all the UI text need to be changed as well because the static text are written
        in default language.
        If request requires the default language the the system returns the static text to improve performance.
    \item Refactoring - When refactoring the static texts the system creates new database records.
        The system can not determine unused keys, and so they are left in database showing in administration dashboard.
    \item Ambiguity - Each key requires translation group to resolve ambiguous keys.
        The group-key pair should always correspond to only one static text.
\end{itemize}

Example usage of the translation system:
\begin{verbatim}
    // MyUiComponent.php
    class MyUiComponent implementes View {
        // UI rendering code ...

        // Provides shorthand methods for static text translation
        use LexiconUnit;

        public function __construct() {
            $this->setLexiconGroup("example.group");
        }
    }

    // MyUiComponent.phtml
    <button>
        <!-- Creates a new pairing: (exmaple.group and Submit) -->
        <!-- and searches for correct translation -->
        <?= $this->tr("Submit") ?>
    <button>
\end{verbatim}



\subsubsection*{Administration Interface}
The administration system automates CRUD operations based on database models.
While this reduces development effort, it introduces constraints:

\begin{itemize}
    \item Workflows - The administration dashboard handles simple CRUD well, but they lack data validation logic.
        If validation is required, then editor behavior needs to implemented.
        The editor behavior is solution to most common create/update feature requests.
        It supports custom form creation and on submit validation while editor takes care of form submission and model creating/updating.
    \item Database structure - It is limitation of the current system that it only works with built-in SQL connections.
        The problem is acknowledged but not addressed further.
    \item Grid - The current system supports custom grid rows definitions, but they have to be descendant of Model class.
        This means, if the row is composed of multiple database tables either each rows explodes into multiple database calls
        during rendering or new abstract model needs to be implemented to wrap results of complicated select queries.
    \item Menu routes - The administration dashboard required side menu for navigation.
        Naive implementation of showing the route segment to the user would work, but it is not ideal for the user,
        since url paths do not support spaces and other special character.
        This problem was solved by implementing proper graph collection described in the Routing section.
\end{itemize}



\subsubsection*{Testing}
The testing component is designed for simplicity, but this simplicity leads to:

\begin{itemize}
    \item Complex testing scenarios - The current system does not support testing beyond simple unit tests.
    \item Implementation - The library used for testing is implemented naively and lacks in implementation quality.
    \item Manual inspection - The results of tests needs to be viewed on the administration dashboard and currently
        does not support running tests automatically and notifying about the results.
\end{itemize}

Example of simple unit tests using the library:
\begin{verbatim}
    // utils.php
    Sptf::test("match arrays", function () {
        Sptf::expect(Arrays::equal([], []))->toBe(true);
        Sptf::expect(Arrays::equal([1], [1]))->toBe(true);
        Sptf::expect(Arrays::equal([1], []))->toBe(false);
        Sptf::expect(Arrays::equal(["foo" => "bar"], ["foo" => "bar"]))
            ->toBe(true);
        Sptf::expect(Arrays::equal(["foo" => "bar"], ["bar" => "foo"]))
            ->toBe(false);
    });
    // Results: PASS 0.00s "match arrays" 5 assertions

    Sptf::test("reverse array", function () {
        $array = [0, 1, 2, 3, 4, 5, 6, 7, 8, 9];

        Sptf::expect([...Arrays::reversed($array)])
            ->compare(fn($a, $b) => Arrays::equal($a, $b))
            ->toBe(array_reverse($array));
    });
    // Results: PASS 0.00s "reverse array" 1 assertions
\end{verbatim}


\subsection{Data Flow and Execution Model}\label{subsec:architecture-data-flow}

Although the design suggests a clean pipeline, the actual execution model includes
additional complexity:

\begin{itemize}
    \item Implicit dependencies between components
    \item Shared mutable state through database connections
    \item Conditional execution paths based on runtime context
\end{itemize}

These factors make the system harder to reason about compared to the idealized
pipeline model.

\subsection{Cross-Cutting Concerns}\label{subsec:cross-cutting-concerns}

Several concerns affect multiple components simultaneously:

\begin{itemize}
    \item \textbf{Performance:} Abstractions introduce overhead, especially in routing
    and database access.
    \item \textbf{Consistency:} Shared data access increases risk of inconsistent state.
    \item \textbf{Scalability:} Tight coupling through shared resources limits horizontal scaling.
    \item \textbf{Debugging:} Indirect interactions between components complicate tracing errors.
\end{itemize}

These concerns highlight the gap between theoretical design and practical implementation.

\subsection{Architectural Trade-offs}\label{subsec:architectural-tradeoffs}

The implementation reflects several key trade-offs:

\begin{itemize}
    \item Simplicity vs. flexibility
    \item Abstraction vs. performance
    \item Automation vs. control
\end{itemize}

In many cases, decisions favor simplicity and rapid development over long-term
scalability and robustness, which is consistent with the prototype nature of the system.

\subsection{Summary}\label{subsec:architecture-summary}

The system architecture exposes the complexity hidden behind the high-level design.
While the design goals emphasize modularity and extensibility, the implementation
reveals constraints imposed by shared data, abstraction layers, and component interaction.

The resulting architecture represents a compromise between theoretical design
principles and practical feasibility, highlighting areas for future refinement.